Design proposal, 23 September 2026. \textbf{SSL} is the project's
working name for a stigmatic singlet lens formed by two Cartesian
surfaces. This document records the intended software and research
contract, not an implemented API or a new experimental result. The
accompanying \href{../../docs/ssl-generator-theory.md}{theory sketch}
connects it to the general continuum formulation in
\href{../../docs/theory/README.md}{the canonical manuscript}.

\section{The objective}\label{the-objective}

Given a fully specified SSL and a declared acoustic apparatus, find a
bounded set of idealized source amplitudes and phases whose computed
acoustic field supports that lens at equilibrium. Return a reproducible
verification report, including unsuccessful outcomes. Do not change the
requested optical system to obtain a more attractive result.

The immediate product is a \textbf{stationary lens compiler}:

\begin{verbatim}
Fixed lens prescription + fixed machine model
                    ↓
Optical construction and apparatus compatibility
                    ↓
Required full-interface mechanical loading
                    ↓
Source amplitudes and phases
                    ↓
Held-command coupled equilibrium and optical verification
                    ↓
Versioned lens state with explicit evidence status
\end{verbatim}

A later controller can form these states or move between them. A
successful formation movie is not a prerequisite for a verified
stationary solution. Conversely, stationary balance alone is not a
formation or stability result. Idealized emitters are an intentional
model choice; electrical calibration is not a version-1 acceptance
requirement.

\section{The user fixes the optical prescription}\label{the-user-fixes-the-optical-prescription}

The version-1 input is

\begin{verbatim}
(n0, n1, n2, z0, z1, z2, front_vertex_z, central_thickness, frame_radius)
\end{verbatim}

together with the design wavelength, clear aperture, branch convention
and accuracy requirements. Distances in files and solvers are metres.

{\def\LTcaptype{none} % do not increment counter
\begin{longtable}[]{@{}
  >{\raggedright\arraybackslash}p{(\linewidth - 2\tabcolsep) * \real{0.5000}}
  >{\raggedright\arraybackslash}p{(\linewidth - 2\tabcolsep) * \real{0.5000}}@{}}
\toprule\noalign{}
\begin{minipage}[b]{\linewidth}\raggedright
Input
\end{minipage} & \begin{minipage}[b]{\linewidth}\raggedright
Meaning and convention
\end{minipage} \\
\midrule\noalign{}
\endhead
\bottomrule\noalign{}
\endlastfoot
\texttt{z0\_m}, \texttt{z1\_m}, \texttt{z2\_m} & Object, shared
intermediate conjugate, and final image, all in one laboratory
coordinate system \\
\texttt{front\_vertex\_z\_m} & Axial position of the first surface
vertex; it can be set to zero by an explicit coordinate convention \\
\texttt{central\_thickness\_m} & Positive vertex-to-vertex separation,
not the separation of the two rim attachment planes \\
\texttt{frame\_radius\_m} & Bounding radius R at which the complete
interfaces meet their supports \\
\texttt{clear\_radius\_m} & Qualified optical radius a, with 0
\textless{} $a\le R$; it is not silently reduced after ray tracing \\
\texttt{indices} & Three homogeneous refractive indices at the stated
optical wavelength and material state \\
\texttt{wavelength\_m} & Design wavelength; specifying it does not imply
a diffraction calculation \\
\texttt{surface\_domain} & Full-frame Cartesian surfaces or explicitly
completed Cartesian pupils \\
\end{longtable}
}

The second vertex is
\texttt{front\_vertex\_z\_m\ +\ central\_thickness\_m}. Thus thickness
must be included when translating the common intermediate point into the
second surface's local coordinates. The user supplies z1; the compiler
does not optimize it. It also leaves the conjugates, vertices,
thickness, aperture and indices unchanged.

Special inputs such as collimated incidence need an explicit
representation and analytic limiting branch, rather than an extremely
large arbitrary distance. Equal indices, a conjugate at a vertex, and
graph-folding branches must be diagnosed separately. The first
implementation may reject unsupported cases with a precise reason rather
than promise every real/virtual geometry.

\subsection{What ``the surfaces are known'' means}\label{what-the-surfaces-are-known-means}

For \textbf{full-frame Cartesian surfaces}, the supplied prescription
selects two regular, physically transmitting branches over $0\le r\le R$.
Their entire shapes, edge heights and enclosed lens volume follow from
that prescription.

For \textbf{Cartesian pupils with non-optical annuli}, the prescription
determines only $0\le r\le a$. The annular profiles on $a<r\le R$
must also be provided or generated by an explicit, deterministic
completion rule with its mounting, smoothness and volume data. They
cannot be ignored in the mechanics or wave problem. Version 1 does not
secretly optimize those annuli.

\textbf{Proposed default, requiring confirmation before API
implementation:} use full-frame Cartesian surfaces, with a separately
declared clear aperture. Retain pupil-plus-annulus mode as an explicit
option because existing runs use it. The two modes are not
interchangeable prescriptions.

\section{The machine is a separate input}\label{the-machine-is-a-separate-input}

The machine model declares the cylinder and end-cap positions, source
support and spatial patterns, active/inactive regions, optical access,
acoustic boundary admittances, phase materials, gravity, attachment
conditions, and source-command limits. Frequency is fixed for a
version-1 compilation; an optional frequency study consists of
separately recorded cases.

Material data must distinguish supplied values, engineering assumptions
and measured quantities. A stationary inviscid acoustic calculation does
not use viscosity merely because viscosity is present in a material
file. Immiscibility, constant surface tension and perfect pinning are
also declared assumptions.

Two operating modes must remain distinct:

\begin{itemize}
\tightlist
\item
  \textbf{Design a machine/fill for one lens:} derive its required rim
  heights and volumes, then test a declared compatible apparatus. This
  is not evidence that an existing machine can change to that lens
  without modification.
\item
  \textbf{Compile for a fixed machine:} treat all geometry, phase
  volumes and attachment data as immutable. Reject incompatible
  prescriptions before source optimization. Do not add fluid, move a rim
  or change an acoustic window without declaring a different apparatus.
\end{itemize}

For a sealed reconfigurable machine, volume and rim compatibility are
central. Changing z0, z1, z2 or thickness generally changes the complete
Cartesian profiles and their volume. Therefore the set of SSLs available
to one machine is a constrained family, not automatically an arbitrary
box of optical inputs. An explicitly designed non-optical annulus may
supply accommodation; pumps or moving supports would instead be
additional apparatus components and control variables, outside this
fixed-machine first version.

\section{Responsibilities of the compiler}\label{responsibilities-of-the-compiler}

\subsection{A. Optical construction}\label{a.-optical-construction}

Construct each unsquared constant-optical-path branch through its
prescribed vertex. Check graph regularity, orientation, Snell
transmission, ray interception, minimum thickness and clearance to the
cylinder boundaries. Trace the ideal pair to verify the common-conjugate
construction independently of emitter synthesis. A Cartesian polynomial
root is not sufficient branch verification.

\subsection{B. Required loading}\label{b.-required-loading}

Calculate nonlinear capillary and hydrostatic loading over both complete
interfaces, modulo the constant pressure multipliers imposed by sealed
volumes. Keep the signed density difference on each interface. Store the
geometry, load, units and sign convention together.

\subsection{C. Source inversion}\label{c.-source-inversion}

At the prescribed geometry, construct the acoustic response to the
declared source space. Optimize complex commands under explicit
amplitude and effort constraints, retaining coherent interference and
forces on both whole faces. Multiple starting points and alternative
discretizations are numerical tools, not changes to the optical
prescription. A local optimizer's failure is not proof that a requested
lens is unreachable.

\subsection{D. Held-command verification}\label{d.-held-command-verification}

Freeze the physical command distribution. Solve the actual coupled
acoustic and mechanical equations, then trace the resulting surfaces at
the fixed detector and fixed illumination. A root solver may start near
the target; that is a legitimate equilibrium search, not a physical
trajectory.

Refine acoustic and surface discretizations with the same source
regions, spatial command patterns, amplitudes, phases and frequency.
Increasing the emitter count or re-optimizing produces a new design, not
a convergence test. Report numerical uncertainty relative to the
requested tolerances.

\section{Acceptance without ambiguous success labels}\label{acceptance-without-ambiguous-success-labels}

Use separate evidence fields rather than one ``success'' flag:

{\def\LTcaptype{none} % do not increment counter
\begin{longtable}[]{@{}
  >{\raggedright\arraybackslash}p{(\linewidth - 2\tabcolsep) * \real{0.5000}}
  >{\raggedright\arraybackslash}p{(\linewidth - 2\tabcolsep) * \real{0.5000}}@{}}
\toprule\noalign{}
\begin{minipage}[b]{\linewidth}\raggedright
Evidence field
\end{minipage} & \begin{minipage}[b]{\linewidth}\raggedright
What it means
\end{minipage} \\
\midrule\noalign{}
\endhead
\bottomrule\noalign{}
\endlastfoot
\texttt{optically\_constructed} & The prescribed ideal pair passes
branch and optical checks \\
\texttt{apparatus\_compatible} & Shapes, supports, phase volumes and
access match the declared machine \\
\texttt{source\_candidate\_found} & An inverse solve found admissible
complex commands \\
\texttt{stationary\_root\_converged} & The full discrete coupled force
residual meets its tolerance \\
\texttt{stationary\_optical\_gates\_met} & Both height errors, maximum
spot and ray coverage pass \\
\texttt{fixed\_command\_convergence\_verified} & Independent refinements
support the reported stationary metrics \\
\texttt{stability\_assessed} & A specified perturbation model has
actually been analysed; include its conclusion and scope \\
\texttt{formation\_verified} & A specified startup trajectory reaches
and maintains the accepted operating state \\
\texttt{transition\_verified} & A specified change between operating
states meets its transient requirements \\
\end{longtable}
}

Unknown, failed and not-applicable are distinct. ``Unreachable'' is
reserved for a valid necessary-condition obstruction or other
certificate within the declared model; an exhausted optimization budget
returns ``no solution found.''

Height error is the maximum absolute cycle-mean lab-coordinate departure
from each prescribed surface on its qualified pupil, without fitted
piston, tilt or focus. Optical error is maximum geometric radius about
the prescribed axis at the fixed detector, with a separate full-coverage
requirement. Neither an RMS radius nor the radius of a clipped surviving
subset substitutes for these gates. Sampling resolution and any
non-certified continuous maximum must be stated.

The intended scales are 10 nm height and approximately 0.5 µm spot.
Executable requests need exact numbers and inequalities: a strict 0.500
µm bound would \textbf{not} accept the saved 0.513 µm result. User
acceptance of an approximate spot does not silently relax the
independent height or coverage criteria.

\section{Proposed software interfaces and records}\label{proposed-software-interfaces-and-records}

The following names are a proposal, not existing callable functions:

\begin{verbatim}
target = construct_ssl(prescription)
compatibility = check_machine(target, machine)
load = required_load(target, machine)
candidate = synthesize_sources(target, load, machine, inverse_settings)
equilibrium = solve_fixed_command(candidate, machine)
evidence = verify_stationary(target, equilibrium, candidate, checks)
\end{verbatim}

Keep immutable prescription, machine, target, source-command, equilibrium
and verification objects separate. Cache keys must include geometry, materials, boundary
operators, source patterns, frequency and numerical settings. A response
matrix for one target cannot silently be reused after changing its
scattering geometry.

Each result belongs in a protected run directory containing requested
and resolved configurations; target and computed surfaces; complex
commands and source-region geometry; optical and force metrics;
source/provenance hashes; mesh and solver histories; and a
human-readable report. A stationary state has no physical time axis. A
trajectory adds actual times, command history, fluid states and
per-frame optical checks.

The notebook is a reader and verification presentation, not an
independent source of physical states. Display both unforced and driven
surfaces, source geometry, harmonic pressure and mean traction
separately, and the full spot diagram. Formation and carrier-wave
animations have different time scales and must never be spliced into one
implied successful experiment.

\section{Extension to an adaptive SSL machine}\label{extension-to-an-adaptive-ssl-machine}

Once stationary states are independently verified, a lens library maps
valid prescriptions to commands and evidence. The first adaptive mode
can be \textbf{select, settle, then image}. A later mode tracks a
time-dependent optical request with bounded transient error.

Endpoint existence does not establish a connecting path. Changes must
preserve the apparatus constraints, respect actuator bandwidth, and
avoid branch loss, interface contact and unacceptable tracking error.
Command interpolation is at most an initial guess for a transition
calculation. Stability and formation may require feedback; specify the
observations and controller rather than assuming exact knowledge of the
liquid surface.

Axial stigmatism is defined for the selected conjugate pair and optical
indices. It does not automatically provide aplanatism, broadband
correction, a large corrected field or a zero diffraction spot. UV
curing would create a different manufacturing workflow, with evolving
shape and refractive index; it is not part of continuously
reconfigurable liquid operation.

\section{What exists and what should be implemented next}\label{what-exists-and-what-should-be-implemented-next}

Existing code supplies Cartesian targets, idealized source regions,
quadratic force synthesis, held-command equilibrium checks, ray tracing
and provenance. The unified fixed-prescription API and its compatibility
contract are still proposed. In particular, the current two-face surface
basis uses fixed rim levels and fixed volume relative to reference
planes. It must not be treated as a general full-frame Cartesian-volume
implementation without extension.

The
\href{../../docs/../artifacts/notebooks/09_emitter_driven_equilibrium.executed.ipynb}{current
stationary notebook} records a 0.513 µm surviving-ray spot, 20.2/24.4 nm
height errors, and one failed ray out of 4001; fixed-command refinement
is unfinished. The
\href{../../docs/../artifacts/notebooks/10_emitter_formation.executed.ipynb}{formation
notebook} is a separate unsuccessful, partial Stokes diagnostic. Neither
demonstrates the proposed adaptive family.

The next implementation milestones are:

\begin{enumerate}
\def\labelenumi{\arabic{enumi}.}
\tightlist
\item
  Fix the prescription schema, coordinate conventions and explicit
  full-frame versus pupil-completion mode; test optical and volume
  compatibility.
\item
  Connect the existing solvers through immutable inputs and explicit
  evidence fields, without changing prescribed geometry during source
  fitting.
\item
  Verify a small, independently selected family with fixed-command
  refinement, including unsuccessful or incompatible cases, before
  mapping a wider range.
\item
  Add stability-qualified operation and then formation/transition
  control.
\end{enumerate}

The research claim to pursue is a reproducible map from
\textbf{specified compatible SSLs to verified acoustic commands}, with
quantified errors, effort and limitations---not universal realizability
inferred from a single lens.

\section{Scientific grounding}\label{scientific-grounding}

The optical construction follows Silva-Lora and Torres, {Superconical
aplanatic ovoid singlet lenses} \cite{silva2020} (2020). Its two-surface
construction is an optical foundation, not evidence that this acoustic
machine realizes the surfaces. The acoustic/mechanical inverse and its
assumptions are developed separately in the
\href{../../docs/ssl-generator-theory.md}{theory sketch} and canonical
manuscript.
